% Unnumbered teaching material; technical statement numbering is preserved.
\clearpage
\section*{An extended bridge: from random events to a physical certificate}
\addcontentsline{toc}{section}{Extended bridge: from random events to a physical certificate}
The previous lessons supplied the vocabulary. These eight additional lessons
slow down the transition from that vocabulary to the physical computations.
They answer a practical question: how can a finite calculation establish a
statement about an entire interval of times? They then examine what it means
for a mathematical rate to have a physical interpretation.

\section*{Lesson 33. A rate is not a probability: build the clock first}
\addcontentsline{toc}{section}{33. A rate is not a probability}
\paragraph{Begin with units.}
A probability is a number between zero and one. A rate has units of inverse
time. A rate of three events per second is perfectly meaningful; a probability
of three is not. If an event has constant rate $r$, its probability of occurring
in a very short interval of length $h$ is $rh+o(h)$. The notation $o(h)$ means
that the error divided by $h$ tends to zero as $h$ tends to zero. It does not
mean that the approximation is exact for a one-second interval.

\paragraph{Derive the waiting law.}
Let $S(t)$ be the probability that no event has happened by time $t$.
Assume that, conditional on no event so far, the chance in the next short
interval is still $rh+o(h)$. Then
\[
 S(t+h)=S(t)(1-rh+o(h)),\qquad S'(t)=-rS(t),\qquad S(0)=1.
\]
The function with this derivative and initial value is $S(t)=e^{-rt}$.
Thus the probability of at least one event by time $t$ is $1-e^{-rt}$.
For $r=2$ per second and $t=0.1$ seconds, this is $1-e^{-0.2}$,
approximately $0.1813$. The short-time estimate $rt=0.2$ is close but different.
For $t=1$ second the estimate $rt=2$ is not a probability at all.

\paragraph{Why the mean wait is $1/r$.}
The density of the first event time is $re^{-rt}$. Its mean is
\[
 \int_0^\infty t r e^{-rt}\,dt
 =[-te^{-rt}]_0^\infty+\int_0^\infty e^{-rt}\,dt=\frac1r.
\]
The integration by parts has merely moved the factor $t$ onto a derivative.
There is also a geometric interpretation: the mean waiting time is the area
under the survival curve $S(t)$.

\paragraph{Several possible exits.}
From state $j$, let the exit rates to other states be $a_{ij}$.
Independent exponential clocks compete. The probability that none has rung
by $t$ is the product $\prod_{i\ne j}e^{-a_{ij}t}$, hence the total exit rate
is $r_j=\sum_{i\ne j}a_{ij}$. Conditional on an exit, its destination is $i$
with probability $a_{ij}/r_j$. This produces the diagonal entry
$M_{jj}=-r_j$ and explains why a column sums to zero. The negative diagonal
does not describe a negative probability; it describes probability leaving.

\paragraph{Try it.}
Suppose the rates from one state to two others are $2$ and $3$ per second.
The mean residence time is $1/5$ second, and the next destination probabilities
are $2/5$ and $3/5$. Multiplying all rates by ten divides every waiting time
by ten. It does not change these destination probabilities or the equilibrium
distribution. Consequently, a percentage time advantage alone supplies no
absolute number of seconds without a physical rate calibration.

\section*{Lesson 34. Replace unequal clocks by one clock and harmless pauses}
\addcontentsline{toc}{section}{34. One clock and harmless pauses}
\paragraph{A deliberately simple two-state example.}
Use a column probability vector and let
\[
 M=\begin{pmatrix}-2&1\\2&-1\end{pmatrix},\qquad
 \rho=2,\qquad
 P=I+M/\rho=\begin{pmatrix}0&1/2\\1&1/2\end{pmatrix}.
\]
The first state exits at rate $2$; the second exits at rate $1$.
Instead ring a clock of rate $2$ in either state. At each ring, choose the
next state using the current column of $P$. In state two, half the rings
leave the state unchanged. Its actual transition rate is therefore
$2\times(1/2)=1$, exactly as required. The pauses are real elapsed time.

\paragraph{Generalize only after the example.}
Choose $\rho$ at least as large as every exit rate. Then $P=I+M/\rho$
has nonnegative entries and unit column sums. Such a matrix is called
stochastic: it sends probability vectors to probability vectors. The identity
$M=\rho(P-I)$ gives
\[
 e^{Mt}=e^{\rho t(P-I)}
 =e^{-\rho t}\sum_{k=0}^\infty\frac{(\rho t)^k}{k!}P^k.
\]
The factorization is allowed because $P$ commutes with $I$.
It would not be legitimate to split an arbitrary matrix exponential into
two factors without checking commutation.

\paragraph{Read every factor as a probability.}
Write $\mu=\rho t$. A Poisson random variable $N$ has probabilities
$\Pr(N=k)=e^{-\mu}\mu^k/k!$. Here $N$ counts clock rings by time $t$.
Conditional on $N=k$, the distribution is $P^kp(0)$. Averaging these
conditional distributions gives the displayed formula. This procedure is
called uniformization. It is an exact continuous-time construction, not a
small-time-step approximation.

At $\mu=1$, the probabilities of zero, one and two rings are respectively
$e^{-1}$, $e^{-1}$ and $e^{-1}/2$. Their sum is about $0.9197$.
The missing $0.0803$ is the probability of three or more rings; it cannot
simply be discarded when a rigorous error bound is required.

\paragraph{Check the two-state equilibrium.}
Solving $M\pi=0$ and $\pi_1+\pi_2=1$ gives $\pi=(1/3,2/3)^\top$.
Starting from $(1,0)^\top$, the exact first component is
$p_1(t)=1/3+(2/3)e^{-3t}$ and the second is its complement.
Differentiating verifies $p'_1=1-3p_1$, which is also the first row of
$Mp$. This small solvable model is a useful check of the matrix convention,
the clock interpretation and any implementation of uniformization.

\section*{Lesson 35. Bound the part of an infinite sum that you did not compute}
\addcontentsline{toc}{section}{35. Bound an infinite remainder}
\paragraph{First learn the geometric comparison.}
For $0\leq q<1$, multiplying $1+q+q^2+\cdots$ by $1-q$ cancels every
term except the first, giving $1/(1-q)$. If a nonnegative series starts with
$a$ and every later term is at most $q$ times its predecessor, its whole sum
is at most $a/(1-q)$. Equality is unnecessary; an upper bound suffices.

For Poisson weights $w_k=e^{-\mu}\mu^k/k!$,
\[
 \frac{w_{k+1}}{w_k}=\frac{\mu}{k+1}.
\]
After keeping terms $0$ through $K$, the first omitted weight is $w_{K+1}$.
Every subsequent ratio is at most $q=\mu/(K+2)$. Provided $q<1$,
\[
 R_K=\sum_{k=K+1}^\infty w_k
 \leq\frac{w_{K+1}}{1-\mu/(K+2)}.
\]
The $K+2$ in the denominator is intentional: it bounds the ratio from the
first omitted term to the next omitted term. Confusing it with the first
omitted term's index is an easy off-by-one mistake.

\paragraph{Work a numerical example.}
Take $\mu=1$ and keep $K=2$. Then $w_3=e^{-1}/6$ and $q=1/4$, so
$R_2\leq 2e^{-1}/9$, approximately $0.08175$. The actual missing mass
is approximately $0.08030$, below the bound. The bound is a little wider
than necessary, but its validity is transparent.

\paragraph{Turn a scalar tail into a vector guarantee.}
Every component of $P^kp(0)$ lies in $[0,1]$. If $s_i$ is the computed
partial sum for component $i$, then $p_i(t)\in[s_i,s_i+R_K]$.
These independent intervals need not sum to exactly one; they are an
enclosure of a probability vector, not themselves a new probability vector.
Renormalizing their endpoints without a proof could invalidate the enclosure.

\paragraph{Include arithmetic error as well as truncation error.}
An exact tail formula does not excuse rounding errors in the partial sum.
The nominal checker computes interval bounds for inputs, exponentials,
matrix products and logarithms. If the entire probability interval is
positive, it can enclose each term $p_i\log(p_i/\pi_i)$ and hence KL.
An interval crossing zero requires a valid alternative treatment or a
failure to certify; a logarithm of an invalid endpoint is not evidence.

In the replayed nominal calculation, at most 281 terms were needed and the
largest reported tail upper bound was below $10^{-45}$. The small tail is
one part of the certificate. The validity of $P$, the input correspondence,
the outward arithmetic and the later comparison inequalities remain
separate necessary parts. More printed digits alone establish none of them.

\section*{Lesson 36. Prove a whole time interval using its endpoints}
\addcontentsline{toc}{section}{36. From endpoints to continuous time}
\paragraph{A descending staircase picture.}
Suppose two distances $D_c(t)$ and $D_h(t)$ are nonincreasing in time.
On a cell $[a,b]$, the cold distance is at least its value at the right
endpoint, while the hot distance is at most its value at the left endpoint:
\[
 D_c(t)\geq D_c(b),\qquad D_h(t)\leq D_h(a).
\]
Therefore a single strict inequality
\[
 D_h(a)<0.9D_c(b)
\]
proves $D_h(t)<0.9D_c(t)$ for every time in that cell.
The endpoints are deliberately mismatched to use the worst case for each
curve. Merely checking $D_h(a)<0.9D_c(a)$ and the analogous inequality at
$b$ would not control a crossing between them.

\paragraph{Why KL decreases here.}
For a stochastic transition matrix $K$ and positive equilibrium $\pi$,
let $r_i=p_i/\pi_i$. When $K\pi=\pi$, the output likelihood ratio is
\[
 \frac{(Kp)_j}{\pi_j}=\sum_i\frac{K_{ji}\pi_i}{\pi_j}r_i.
\]
The coefficients in this sum are nonnegative and sum to one. Convexity of
$f(r)=r\log r$ says that $f$ of this weighted average is no greater than
the weighted average of $f$. Multiply by $\pi_j$, sum over $j$, and use
$\sum_jK_{ji}=1$ to obtain
\[
 D(Kp\Vert\pi)\leq D(p\Vert\pi).
\]
Apply this to $K=e^{M(b-a)}$. This argument proves the monotonicity used
above. It does not say that $D_h(t)/D_c(t)$ is monotone.

\paragraph{A hand-checkable cell.}
Suppose outward bounds give $D_h(a)\leq0.0080$ and
$D_c(b)\geq0.0090$. Then $0.9D_c(b)-D_h(a)\geq0.0001>0$.
Every interior time passes the ten-percent comparison. If the computed
margin interval were $[-0.0001,0.0002]$, no conclusion would follow;
a narrower cell might resolve it.

\paragraph{Join the cells and check the floor.}
A finite list $t_0<t_1<\cdots<t_N$ covers the entire closed interval
$[t_0,t_N]$ without gaps. Proving the endpoint inequality on every cell
proves the full window. Since cold KL decreases, a lower bound
$D_c(t_N)\geq10^{-3}$ also proves that floor throughout the window.
The nominal replay used 80 cells and a minimum certified comparison margin
greater than $0.0004408$. Its fixed endpoints were approximately
$58.89269$ and $117.78539$ in attempt-time units. The script separately
encloses the slow time to justify the window-length requirement; a decimal
time chosen from a numerical eigenvalue does not become exact merely by
being described as a multiple of that numerical estimate.

\section*{Lesson 37. Prove a whole band of target distances}
\addcontentsline{toc}{section}{37. A whole band of target distances}
\paragraph{Reverse the graph in your mind.}
For a positive target $\delta$ below the initial distance, define
$t_x(\delta)$ as the first time preparation $x$ reaches that target.
A smaller target means getting closer to equilibrium, so it takes longer.
Thus if $\delta_a<\delta_b$, then $t_x(\delta_a)\geq t_x(\delta_b)$.
For the irreducible reversible model and a nonequilibrium preparation,
KL strictly decreases until its limiting equilibrium value; positive targets
below the initial value have unique crossings.

\paragraph{A root estimate proposes; a sign bracket verifies.}
A numerical solver might propose $t=10.1234$. To verify a bracket $[L,U]$,
enclose $D(L)-\delta$ strictly above zero and $D(U)-\delta$ strictly below
zero. Continuity gives a crossing between them; strict decrease gives
uniqueness. An approximate residual close to zero is weaker than these
two certified signs. We also need $0<L<U$ for the ratio bound below.

\paragraph{Choose the unfavorable corners.}
For every $\delta\in[\delta_a,\delta_b]$,
\[
 t_h(\delta)\leq t_h(\delta_a)\leq U_h(\delta_a),\qquad
 t_c(\delta)\geq t_c(\delta_b)\geq L_c(\delta_b)>0.
\]
Division by a positive denominator preserves the intended order, giving
\[
 1-\frac{t_h(\delta)}{t_c(\delta)}
 \geq 1-\frac{U_h(\delta_a)}{L_c(\delta_b)}.
\]
If the right side exceeds $0.05$, every target in that cell has at least
five-percent time saving. Notice that the numerator and denominator use
opposite ends of the target cell. This conservative mismatch is exactly
what covers its interior.

\paragraph{Try the ratio with ordinary numbers.}
Suppose the hottest trajectory's slowest permitted arrival is at most
$90$ seconds and the cold trajectory's fastest permitted arrival is at
least $100$ seconds. The saving is at least $1-90/100=10\%$.
If the denominator were instead only known to exceed $92$, the bound would
fall to about $2.17\%$, even if all sampled same-target comparisons looked
better. The proof must use the bounds it has actually established.

The nominal certificate covers $[10^{-3},10^{-2}]$ with 60 target cells.
Its conservative lower bound on time saving is approximately $8.12245\%$.
This is a guaranteed lower bound for that nominal model, not the exact
minimum. The later eight-obligation certificate handles the full uncertain
four-state parameter box. It is a distinct proof, not a conclusion obtained
by silently extrapolating the nominal calculation.

\section*{Lesson 38. Derive a transition rate from an energy exchange}
\addcontentsline{toc}{section}{38. From energy exchange to transition rate}
\paragraph{Begin with one uphill move.}
Let two system energies be $E_j=1$ and $E_i=3$, with barrier energy $B=5$.
Draw a nonnegative reservoir energy $Y$ with density $\beta e^{-\beta Y}$.
The move from $j$ is allowed when $1+Y\geq5$, namely when $Y\geq4$.
At acceptance, set $Y'=Y+1-3=Y-2$. It is nonnegative, and
$E_i+Y'=3+Y-2=1+Y=E_j+Y$: combined energy is conserved.
The reverse contact sees the same total energy and can undo the move.

\paragraph{Integrate the acceptance probability.}
For any $c\geq0$,
\[
 \Pr(Y\geq c)=\int_c^\infty\beta e^{-\beta y}\,dy=e^{-\beta c}.
\]
If contacts for this pair arrive at rate $\nu_0$, the accepted rate from
$j$ to $i$ is therefore
\[
 a_{ij}=\nu_0e^{-\beta(B-E_j)}.
\]
The reverse rate is $a_{ji}=\nu_0e^{-\beta(B-E_i)}$.
Their ratio is $e^{-\beta(E_i-E_j)}$. Multiplying by Gibbs weights
gives the stronger flux identity
\[
 a_{ij}e^{-\beta E_j}=a_{ji}e^{-\beta E_i}
 =\nu_0e^{-\beta B},
\]
which is detailed balance before the common normalization is inserted.

\paragraph{Separate an always-required gate from the uphill cost.}
Write $h=B-\max(E_i,E_j)$. Since
$B-E_j=h+\max(E_i-E_j,0)$, the rate becomes
\[
 a_{ij}=\nu_0e^{-\beta h}\min\{1,e^{-\beta(E_i-E_j)}\}.
\]
In the example at $\beta=1$, the common gate contributes $e^{-2}$,
the uphill acceptance contributes another $e^{-2}$, and their product is
$e^{-4}$. The reverse move has only the common $e^{-2}$ gate.

\paragraph{Why simulation may omit attempts but not their time.}
One can retain only contacts with $Y\geq h$. Their clock has rate
$\nu_0e^{-\beta h}$. Conditional on passing the gate,
\[
 \Pr(Y-h\geq z\mid Y\geq h)
 =\frac{e^{-\beta(h+z)}}{e^{-\beta h}}=e^{-\beta z}.
\]
Thus a retained mark is $h$ plus a fresh exponential energy. This proves
the thinning rule used in the event simulation. Counting every attempted
contact as an accepted transition would alter the process; ignoring
rejected contacts without adjusting the clock would alter physical time.

\paragraph{What this construction establishes.}
Under independent fresh thermal marks, independent Poisson pair clocks,
constant reservoir density of energy states, and instantaneous updates,
the rate identity holds exactly for every allowed barrier and energy.
It supplies a stochastic bath interpretation of the same four-state model.
It does not identify a material or derive an autonomous microscopic bath.
The replayed simulation of 500,000 trajectories per preparation agrees
with direct evolution at five times to within $0.00118$ per component.
That is a diagnostic check of the simulation; the rate identity above is
the reason the mathematical generators coincide.

\section*{Lesson 39. An intermediate state introduces a second clock}
\addcontentsline{toc}{section}{39. Intermediate states and two clocks}
\paragraph{Insert a waiting room between two wells.}
Suppose the desired well-to-well rates are $u$ from $j$ to $i$ and $v$
from $i$ to $j$. Insert state $s$. Use entrance rates $2u$ from $j$
and $2v$ from $i$, and exit rate $b$ from $s$ to each well. On entering
$s$, the next well is selected with probabilities $1/2$ and $1/2$.
Half the visits return to the well just left.

If the clock is stopped during every visit to $s$, the successful transition
rate from $j$ to $i$ is $2u\times(1/2)=u$. This is the trace clock:
only time actually spent in the wells is retained. With ordinary physical
time, each visit adds a mean wait $1/(2b)$. Those waits cannot be removed
when claiming a physical completion time.

\paragraph{Derive the block identity one entry at a time.}
Order the wells as $(j,i)$. Their direct block, the entrance block, the
exit block and the intermediate block are
\[
 Q_{ww}=\begin{pmatrix}-2u&0\\0&-2v\end{pmatrix},\quad
 Q_{sw}=\begin{pmatrix}2u&2v\end{pmatrix},\quad
 Q_{ws}=\begin{pmatrix}b\\b\end{pmatrix},\quad Q_{ss}=-2b.
\]
The subscripts give destination first, as everywhere else. Multiplication
now gives
\[
 Q_{ww}-Q_{ws}Q_{ss}^{-1}Q_{sw}
 =\begin{pmatrix}-2u&0\\0&-2v\end{pmatrix}
  +\begin{pmatrix}u&v\\u&v\end{pmatrix}
 =\begin{pmatrix}-u&v\\u&-v\end{pmatrix}.
\]
This is exactly the intended two-well generator. Adding the six pair
contributions gives the four-well identity in the technical report.
This exact algebra identifies the trace generator; it does not assert that
the physical-time well probabilities solve that same four-state equation.

\paragraph{Equilibrium and preparation also change.}
In the expanded model, each intermediate has energy $B$ and phase-space
weight $\eta>0$. Its equilibrium weight is $\eta e^{-B}$, compared
with $e^{-E_j}$ for a well. Taking $b=2/\eta$ in attempt-time units
gives mean intermediate residence time $\eta/4$ and matched equilibrium
fluxes. At preparation temperature $T$, all ten states must be prepared
with weights $g_a e^{-E_a/T}$. Preparing only the wells would be a different
experiment. The phase-space weight $\eta$ need not be an integer count.

At $\eta=0.05$, the nominal cold and hot preparations place about $1.90\%$
and $4.07\%$ of their mass in intermediates. Full ten-state calculations
retain those populations and every residence time. Their finite-time
advantage persists at the tested points. The expanded system also has
higher energies and faster rates outside the original all-state caps.
It is a separate test of a modeling assumption, not a new certificate
inside the original constrained four-state class.

\section*{Lesson 40. Read a computational result at the right strength}
\addcontentsline{toc}{section}{40. Read the evidence at the right strength}
\paragraph{Four questions before accepting a result.}
What object was calculated? What inputs were used? Which errors were bounded?
Which variables were covered continuously? An answer such as ``70 digits''
only addresses numerical precision, and even then does not itself establish
a bound on every implementation or modeling error.

The independent ten-state calculation directly constructs and exponentiates
the full generator at 70-digit working precision. It agrees with the fast
spectral implementation within $6.55\times10^{-13}$ in KL at five tested
times. Different formulas agreeing is useful evidence against a numerical
implementation mistake. Neither implementation uses outward error bounds
for that comparison, so it is a high-precision diagnostic, not an interval
certificate.

\paragraph{Interpret a large sample precisely.}
For each of three fixed $\eta$ values, the study tests 256 corners and 500
seeded interior parameter points. The total is $3(256+500)=2268$.
Each point is further tested on finite target and time grids. No sampled
case fails the declared benchmarks. This says exactly what happened on
those grids. It does not exclude a narrow failure between points, between
time samples, between targets, or at an untested value of $\eta$.
The continuous four-state certificate cannot be transferred to this
ten-state extension merely because their trace generators agree.

\paragraph{A percentage needs a denominator and a clock.}
At $\eta=0.05$ and target $10^{-3}$, the independently computed arrival
times are approximately $210.67249$ for cold and $190.05131$ for hot.
Using the cold time as denominator gives
\[
 \frac{210.67249-190.05131}{210.67249}\approx0.0978826,
\]
or about $9.78826\%$. Dividing by the hot time would answer a different
percentage question. These are attempt-time units, not calibrated seconds.
The calculation includes post-quench intermediate residence times but does
not include the cost or duration of preparing the hotter initial state.

\paragraph{Know what a program's success message establishes.}
A zero exit code means that the process completed without a reported
execution failure. Some exploratory programs record failures in JSON
without returning a nonzero code. Read the recorded failure counts, input
hashes, unresolved cells, parameter scope and numerical residuals. A
certificate runner must also reject incomplete coverage. A numerical
search that found nothing has a different meaning from a complete proof
that nothing exists.

\paragraph{The final handoff to advanced mathematics.}
You can now read the physical section by following a chain of explicit
claims: construct valid rates; establish equilibrium; evolve preparations;
bound numerical error; prove continuous comparisons; and interpret the
clock and physical assumptions. The advanced spectral and interval
methods make that chain efficient over an uncertain parameter box.
They do not replace any link. When a later formula looks forbidding,
identify which link it serves before trying to manipulate its symbols.
